\section{Locator implementation}
\label{sec:pg-storage-sequence-indirection}

Sequence locator are implemented as 40-byte struct in C++. The struct must also include the Postgres header so it's effectively 56 bytes in total.
The Postgres header stores alphabet code and base storage unit as \texttt{TYPMOD} field (custom field for storing type modifier).
I don't provide small-sequence optimisation in BioseqDB, because that would complicate the implementation (instead of storing locator for small sequences we can potentially 
store sequence directly in the user table. This would require storing it's padding so potential maximum sequence length that could be stored inline would be $(40B - 1B)*8/2 = 156$ symbols for DNA4).

\noindent\begin{minipage}{\linewidth}
\begin{lstlisting}[style=bioseq-cpp26,numbers=none,
  basicstyle=\footnotesize\ttfamily,
  caption={The 40-byte locator stored after the 16-byte sequence header.},
  label={lst:storage-sequence-locator}]
struct sequence_locator {
  std::uint8_t  payload_version;   // currently 3
  std::uint8_t  bucket_index;
  std::uint16_t reserved_16;
  std::uint32_t shard_index;
  std::uint32_t fragment_count;
  std::uint32_t reserved_32;
  std::uint64_t logical_length;
  std::uint64_t first_fragment_row;
  std::uint64_t source_epoch;
};
\end{lstlisting}
\end{minipage}


\begin{itemize}
  \item the length bucket and shard that hold the sequence,
  \item its first fragment row and number of fragments,
  \item its logical length in symbols, and
  \item its \textbf{source epoch}, an identifier for this immutable
    stored sequence version.
\end{itemize}

The current locator format is version 3. Readers check the version, non-zero
epoch, fragment count, and final padding before accepting it. They also check
the full registered source identity, which includes the parent table and
column. A matching fragment row number by itself is insufficient: another
relation or source version could use the same number.

\subsection{Symbols, packed units, fragments, and pages}

The units have different meanings. A packed storage unit of $B$ bytes holds
$L$ symbols according to the alphabet encoding from
Section~\ref{sec:compact-sequence-range}. A fragment reserves 1,024 of these
units. Thus a full fragment has
\[
  \text{payload bytes}=1024B,
  \qquad
  \text{symbol capacity}=1024L.
\]
For DNA4 and 32-bit units, $B=4$ and $L=16$: one fragment has 4,096
payload bytes and room for 16,384 symbols. With 64-bit units, it has 8,192
payload bytes and room for 32,768 symbols. A fragment is consequently not a
synonym for a PostgreSQL page, even when their sizes happen to be similar.

Consider a 40,000-symbol DNA4 sequence stored with 32-bit units. It needs
three fragment rows. The last row uses only 7,232 of its symbol positions,
but its packed payload allocation still has the full fragment size.
The locator's length prevents a read from returning that unused capacity.
This is separate from unused symbol positions in the final packed unit.
If the first fragment row number is $r_0$, fragment $i$ is at $r_0+i$.
Consecutive row numbers avoid a separate stored address for every fragment.
The locator needs only the first row number and the fragment count.

\begin{figure}[htbp]
\centering
\begin{tikzpicture}[x=1cm,y=1cm,font=\footnotesize]
  \node[anchor=west,font=\bfseries] at (0,4.75)
    {40,000 DNA4 symbols, packed in 32-bit units};
  \node[anchor=east] at (1.4,3.75) {symbols};
  \draw[draw=black!55,fill=black!4] (1.7,3.35) rectangle (11.9,4.15);
  \foreach \x in {5.1,8.5} \draw[black!55] (\x,3.35)--(\x,4.15);
  \draw[fill=black!20,draw=black!55] (10.0,3.35) rectangle (11.9,4.15);
  \node[align=center] at (3.4,3.75) {fragment\\0};
  \node[align=center] at (6.8,3.75) {fragment\\1};
  \node[align=center] at (9.25,3.75) {fragment\\2};
  \node at (10.95,3.75) {unused};
  \foreach \x/\n in {1.7/0,5.1/{16\,384},8.5/{32\,768},10/{40\,000}}
    \node[above=2pt] at (\x,4.15) {$\n$};
  \draw[storageOrange,line width=3pt] (7.925,2.9)--(8.755,2.9);
  \draw[storageOrange] (7.925,2.75)--(7.925,3.2);
  \draw[storageOrange] (8.755,2.75)--(8.755,3.2);
  \node[anchor=east,text=storageOrange] at (7.65,2.9)
    {request $[30\,000,34\,000)$};
  \node[storage cell,fill=storageBlue!12,minimum width=2.4cm]
    (locator) at (1.75,1.2) {parent locator\\first row $r_0$};
  \foreach \x/\i/\r in {4.55/0/{r_0},7.6/1/{r_0+1},10.65/2/{r_0+2}} {
    \node[storage cell,fill=storageGreen!10,minimum width=2.5cm,
      minimum height=1.15cm] (f\i) at (\x,1.2)
      {row $\r$\\4,096 payload bytes};
  }
  \draw[storage arrow] (locator.east)--(f0.west);
  \draw[storage arrow,storageOrange] (8.1,2.65)--(f1.north);
  \draw[storage arrow,storageOrange] (8.65,2.65)--(f2.north);
  \node[anchor=west,text=black!65] at (.5,.15)
    {The request uses fragments 1 and 2; it does not ask for fragment 0.};
\end{tikzpicture}
\caption{A symbol interval becomes fragment-row requests. Equal-width
fragments reserve equal packed space. Only the used prefix of the final
fragment belongs to the sequence. Underlying column chunks and database
pages can contain additional data, so two requested fragments do not imply
exactly two physical reads.}
\label{fig:sequence-fragment-flow}
\end{figure}

Let $n$ be the sequence length and $F$ the number of symbols per fragment.
For a requested interval $[a,b)$ with $0\leq a<b\leq n$, the required
fragment numbers run from
\[
  \left\lfloor a/F\right\rfloor
  \quad\text{through}\quad
  \left\lfloor (b-1)/F\right\rfloor.
\]
The subtraction by one keeps a request ending exactly at a boundary from
reading the following fragment. This is the physical counterpart of the
half-open slice convention introduced earlier.

\subsection{Length buckets separate short and long values}

A \textbf{length bucket} groups sequences by their logical number of
symbols. A \textbf{shard} is an internal fragment relation within one bucket.
Both belong to a particular sequence column. Length buckets are local
storage groups, not partitions sent to different database servers.
For known-length input, a long reference sequence does not fill the active
shard used for short reads.

The current layout has 35 buckets. The first covers lengths through 256,
the next through 512, and the next through 1,024. The doubling ranges
continue through $2^{41}$. The final bucket covers larger lengths.
Relations are created when needed. A table does not eagerly create every
possible bucket and shard.

\begin{figure}[htbp]
\centering
\begin{tikzpicture}[x=1cm,y=1cm,font=\footnotesize]
  \node[anchor=west,font=\bfseries] at (0,4.75)
    {One sequence column, several length buckets};
  \foreach \x/\b/\range in {1.65/0/{0--256},4.6/1/{257--512},7.55/2/{513--1,024},10.5/34/{over $2^{41}$}} {
    \draw[rounded corners,black!55,fill=black!3]
      (\x-1.25,.8) rectangle (\x+1.25,3.8);
    \node[align=center] at (\x,3.3) {bucket \b\\\range};
    \draw[fill=storageGreen!12,draw=storageGreen]
      (\x-1.02,1.2) rectangle (\x+1.02,2.55);
    \node at (\x,2.22) {one shard};
  }
  \node[storage cell,fill=storageBlue!15,minimum width=1.45cm]
    at (1.65,1.65) {200 symbols};
  \node[storage cell,fill=storageBlue!15,minimum width=1.45cm]
    at (4.6,1.65) {300 symbols};
  \node[storage cell,fill=storageBlue!15,minimum width=1.45cm]
    at (7.55,1.65) {900 symbols};
  \node[storage cell,fill=storageOrange!15,minimum width=1.45cm]
    at (10.5,1.65) {300 symbols};
  \node[anchor=west,text=storageBlue] at (.4,.22)
    {Blue: final length known before writing};
  \node[anchor=west,text=storageOrange] at (.4,-.4)
    {Orange: streaming input; final length was unknown at placement};
  \node[anchor=west,align=left] at (.4,-1.25)
    {Middle buckets are omitted. Every bucket can acquire more shards.\\
     All fragments of one sequence stay together in one shard.};
\end{tikzpicture}
\caption{Length-based placement and its streaming exception. Known-length
inputs choose their length range. Streaming inputs currently use bucket 34,
even when the completed sequence is short.}
\label{fig:storage-length-buckets}
\end{figure}

Known-length insertion can choose a shard before writing:
\begin{enumerate}
  \item Choose the length bucket and calculate the space for its full
    fragment rows, including padding and row headers.
  \item Check the active shard's recorded and pending use against its
    configured size limit.
  \item If the sequence fits but the active shard lacks space, open a
    new shard. If the sequence alone exceeds the limit, reject it.
  \item Reserve consecutive fragment rows in the selected shard, write
    them, and retain their first row number in the parent locator.
\end{enumerate}

The default configured shard size is 1~TiB. This is an allocation limit based
on fragment-row sizes, not a promise about compressed disk size. Streaming
input currently has a narrower check: it verifies the growing sequence's
size, but does not first reserve space for its unknown total size in the
active shard. It also remains in the highest bucket after the final length
becomes known. Thus the limit is not a strict total-size guarantee for a
shard receiving streaming inputs.

Transaction-scoped advisory locks coordinate allocation within a parent
column and bucket, and writes within a selected shard. The compact lock key
only coordinates writers. Locator validation uses the complete registered
identity, so a lock-key collision can cause extra waiting without making
unrelated sequence addresses interchangeable.
